% ---------------------------------------------------

\subsection{Hull}

The hull is a trimaran design with a displacement-type main hull (vaka) and two outrigger hulls (amas) connected by arms. The trimaran configuration provides superior stability---the amas act as counterbalances during heeling, preventing capsize.

\textbf{Fabrication:} The hull is manufactured via fiberglass wet layup in a two-part female mold (bow and stern halves, bolted together). The mold is 3D-printed in standard PLA and surface-finished with Rustoleum filler primer, SprayMax 2K topcoat, PARTALL paste wax, and PVA film release agent. Seven sheets of 4~oz plain-weave fiberglass cloth are laid up with TotalBoat 5:1 marine epoxy resin, yielding approximately 1~mm wall thickness. The keel mounting area receives additional fiberglass layers and an embedded metal plate bonded with resin for structural reinforcement. The stern is similarly reinforced and fitted with a removable mounting plate for rudder installation. The deck will be fabricated from foam separately and bonded to the hull after curing, allowing for easier access to the interior during assembly.

\textbf{Responsible:} Harrison LeTourneau.

\textbf{Completed Tasks:} Hull mold printed and surface-finished; fiberglass layup complete.

\textbf{Pending Tasks:} Mounting hardware installation; deck fabrication and bonding; final sanding; balancing and weight distribution adjustments; static float testing.

% ---------------------------------------------------

\subsection{Sail}

The sail is a self-trimming rigid wingsail consisting of a main sail element and a tail mounted aft on a long rod. The tail provides passive aerodynamic torque to rotate the main sail to the optimal angle of attack relative to the wind, minimizing power consumption compared to motor-actuated designs.

\textbf{Fabrication:} The main sail and tail are carved from XPS rigid foam board and covered in fiberglass for durability. The mast is a 3~ft section of 6061 aluminum round tube, mounted on a bearing at the deck to allow free rotation. A servo housed inside the sail body drives tail rotation via wire linkage, providing the self-trimming function.

\textbf{Responsible:} Charbel Salloum, Connor Stevens.

\textbf{Completed Tasks:} Foam shaping.

\textbf{Pending Tasks:} Fiberglass covering; servo integration; linkage fabrication; mounting hardware installation.

% ---------------------------------------------------

\subsection{Rudder}

The rudder is a spade-type design, chosen for simplicity and ease of modification. It is free-floating (not attached to a skeg), which simplifies construction at the cost of reduced protection from debris impact.

\textbf{Fabrication:} The rudder shaft is a 304 stainless steel round bar (0.25~in. $\times$ 12~in.) with a 3D printed rudder plate that clamps around the bar. It mounts to the stern via a fore/aft clamp mechanism using two rear mounting plates (Aluminium, 0.065~in.) secured with 1/4-20 hex bolts, nylon insert lock nuts, and stainless washers. Ceramic ball bearings (R188, 1/4 $\times$ 1/2 $\times$ 3/16~in.) enable free rotation. A DS3218 MG digital servo drives the rudder via a 25-tooth spline-to-1/4-in. shaft coupler.

\textbf{Responsible:} Corbin Barney, Jared Pratt.

\textbf{Completed Tasks:} Rudder shaft and plate fabricated.

\textbf{Pending Tasks:} Bearing installation; mounting hardware installation; servo integration; linkage fabrication.

% ---------------------------------------------------

\subsection{Keel}

The keel is a fin keel with a weighted bulb, chosen for manufacturability and effective center-of-mass control. The fin provides lateral resistance (generating keel lift to counteract sail lift), while the bulb lowers the center of gravity for improved stability.

\textbf{Fabrication:} The keel frame uses 304 stainless steel---a square bar post (0.5~in. $\times$ 7~in.) with plate mounting brackets. Two mirrored PLA fins are clamped around the square post and secured with 3/16-in. binding posts. The bulb is a 304 SS round tube (2~in. $\times$ 8.5~in.) filled with up to 2~kg of steel balls for ballast. The assembly is wrapped in fiber glass with epoxy resin for structural rigidity.

\textbf{Responsible:} Corbin Barney, Jared Pratt.

\textbf{Completed Tasks:} Keel post and mounting brackets fabricated; fin molds printed.

\textbf{Pending Tasks:} Fin layup; mounting hardware installation.

% ---------------------------------------------------

\subsection{Complete Boat}

Integration of hull, keel, rudder, and sail subsystems requires careful alignment of mounting points and verification of center of effort (COE) and center of lateral resistance (CLR) positioning. The COE (where sail lift acts) should be slightly aft of the CLR (where keel lift acts) to achieve weather helm---a natural tendency to turn into the wind that improves stability and safety. We will also need to ensure the rudder has sufficient authority to counteract the aerodynamic forces on the sail, which may require adjusting the rudder size or position after initial testing. Finally, the vessel will be balanced by adjusting the amount of ballast in the keel bulb and the distribution of weight within the hull to achieve the desired waterline, freeboard, and beam rock characteristics.

\textbf{Testing:} Static float tests will verify waterline, freeboard, and draft. Tow tests will confirm rudder authority and directional stability before autonomous operation. Weight distribution will be adjusted by repositioning steel ballast in the keel bulb.

\textbf{Responsible:} ALAN Team.

% ---------------------------------------------------
